\sezione{2}{Anatomia di un laboratorio}{lab.conf e file .startup}

\subsection{La cartella del laboratorio}
Un laboratorio Kathará è semplicemente una \textbf{cartella}. Il nome dei file è importante: Kathará riconosce i dispositivi dai nomi.
\begin{center}
\begin{tikzpicture}[
  f/.style={anchor=west,font=\ttfamily\small,text=navy},
  d/.style={anchor=west,font=\ttfamily\small\bfseries,text=ambra!70!black},
  c/.style={anchor=west,font=\headfont\footnotesize,text=grigio}]
  \node[d] (root) at (0,0) {kathara-lab\_basic-ipv4/};
  \draw[grigio!60,line width=.8pt] (.25,-0.25) |- (.55,-0.6);
  \node[f] at (.6,-0.6) {lab.conf};
  \node[c] at (4.2,-0.6) {topologia: dispositivi, interfacce, LAN, immagini};
  \draw[grigio!60,line width=.8pt] (.25,-0.85) |- (.55,-1.2);
  \node[f] at (.6,-1.2) {pc1.startup};
  \node[c] at (4.2,-1.2) {comandi eseguiti all'avvio di pc1};
  \draw[grigio!60,line width=.8pt] (.25,-1.45) |- (.55,-1.8);
  \node[f] at (.6,-1.8) {pc2.startup};
  \node[c] at (4.2,-1.8) {comandi eseguiti all'avvio di pc2};
  \draw[grigio!60,line width=.8pt] (.25,-2.05) |- (.55,-2.4);
  \node[f] at (.6,-2.4) {pc3.startup};
  \node[c] at (4.2,-2.4) {comandi eseguiti all'avvio di pc3};
  \draw[grigio!60,line width=.8pt] (.25,-2.65) |- (.55,-3.0);
  \node[f] at (.6,-3.0) {r1.startup};
  \node[c] at (4.2,-3.0) {comandi eseguiti all'avvio di r1};
  \draw[grigio!60,line width=.8pt] (.25,-3.25) |- (.55,-3.6);
  \node[f] at (.6,-3.6) {r2.startup};
  \node[c] at (4.2,-3.6) {comandi eseguiti all'avvio di r2};
  \draw[grigio!60,line width=.8pt] (.25,-3.85) |- (.55,-4.2);
  \node[d] at (.6,-4.2) {shared/};
  \node[c] at (4.2,-4.2) {(facoltativa) cartella condivisa, montata come /shared};
\end{tikzpicture}
\end{center}
\begin{tecnica}[Regole dei nomi]
\begin{itemize}
  \item Un dispositivo esiste se compare nel \texttt{lab.conf} (o se ha una cartella o un file \texttt{.startup} con il suo nome).
  \item \texttt{pc1.startup} è uno script di shell eseguito \textbf{dentro} \texttt{pc1} appena il container parte.
  \item Una eventuale cartella \texttt{pc1/} viene copiata nel file system di \texttt{pc1} (per esempio \texttt{pc1/etc/...} finisce in \texttt{/etc/...}): serve nei lab con servizi come DNS o web server. In questo lab non c'è.
\end{itemize}
\end{tecnica}

\subsection{Il file lab.conf}
La sintassi è una sequenza di assegnamenti \texttt{dispositivo[argomento]=valore}. Le righe che iniziano con \texttt{\#} sono commenti.
\begin{filecfg}{lab.conf}
°# --- router r1: eth0 su A, eth1 su B, MAC scelti a mano ---°
§r1[0]§="A/00:00:00:00:00:a1"
§r1[1]§="B/00:00:00:00:00:b1"
r1[image]="kathara/core"
r1[ipv6]="false"

§r2[0]§="C/00:00:00:00:00:c1"
§r2[1]§="B/00:00:00:00:00:b2"
r2[image]="kathara/core"
r2[ipv6]="false"

§pc1[0]§="A/00:00:00:00:00:01"
pc1[image]="kathara/core"
pc1[ipv6]="false"

§pc2[0]§="C/00:00:00:00:00:02"
pc2[image]="kathara/core"
pc2[ipv6]="false"

§pc3[0]§="C/00:00:00:00:00:03"
pc3[image]="kathara/core"
pc3[ipv6]="false"

°# --- dispositivo speciale per catturare il traffico ---°
wireshark[bridged]=true
wireshark[port]="3000:3000"
wireshark[image]="lscr.io/linuxserver/wireshark"
wireshark[num_terms]=0
\end{filecfg}

\subsubsection{Come si legge una riga di interfaccia}
\begin{center}
\begin{tikzpicture}[part/.style={font=\ttfamily\large,inner sep=2pt},
  ex/.style={font=\headfont\footnotesize,align=center,text=grigio}]
  \node[part,text=navy] (a) at (0,0) {r1};
  \node[part,text=grigio,anchor=west] (b) at (a.east) {[};
  \node[part,text=rosso,anchor=west] (c) at (b.east) {0};
  \node[part,text=grigio,anchor=west] (d) at (c.east) {]="};
  \node[part,text=ambra!80!black,anchor=west] (e) at (d.east) {A};
  \node[part,text=grigio,anchor=west] (f) at (e.east) {/};
  \node[part,text=viola,anchor=west] (g) at (f.east) {00:00:00:00:00:a1};
  \node[part,text=grigio,anchor=west] (h) at (g.east) {"};
  \draw[navy,-{Stealth}] (a.south) -- ++(0,-.6) node[below,ex,text=navy]{dispositivo};
  \draw[rosso,-{Stealth}] (c.south) -- ++(0,-1.25) node[below,ex,text=rosso]{interfaccia\\[-1pt]0 $\to$ eth0};
  \draw[ambra,-{Stealth}] (e.south) -- ++(0,-.6) node[below,ex,text=ambra!70!black]{collision\\[-1pt]domain};
  \draw[viola,-{Stealth}] (g.south) -- ++(0,-.6) node[below,ex,text=viola]{MAC address (facoltativo;\\[-1pt]se manca è casuale)};
\end{tikzpicture}
\end{center}
«Collega l'interfaccia \texttt{eth0} di \texttt{r1} alla LAN \texttt{A} e dagli il MAC \texttt{00:00:00:00:00:a1}.»

\subsubsection{Le opzioni usate nel lab}
\begin{center}
\rowcolors{2}{chiaro}{white}
\begin{tabularx}{\linewidth}{>{\ttfamily}l Y}
\rowcolor{navy}\thd{Opzione} & \thd{Significato}\\
image="kathara/core" & immagine Docker del dispositivo: una Debian con gli strumenti di rete (\texttt{ip}, \texttt{ping}, \texttt{traceroute}, \texttt{arp}, \dots)\\
ipv6="false" & disattiva IPv6: le catture restano pulite, senza i messaggi IPv6 automatici\\
bridged=true & aggiunge un'interfaccia collegata alla rete dell'host tramite NAT (sarà \texttt{eth0})\\
port="3000:3000" & inoltra la porta 3000 dell'host alla porta 3000 del container (\texttt{HOST:GUEST})\\
num\_terms=0 & numero di terminali da aprire per il dispositivo: 0 = nessuno\\
\end{tabularx}
\end{center}
Altre opzioni utili: \texttt{mem} e \texttt{cpus} (limiti di risorse), \texttt{exec} (comando all'avvio), \texttt{sysctl} (parametri del kernel di rete), \texttt{env} (variabili d'ambiente), \texttt{shell} (shell usata da \texttt{connect}). Si possono anche aggiungere metadati come \texttt{LAB\_NAME}, \texttt{LAB\_DESCRIPTION}, \texttt{LAB\_AUTHOR}, mostrati all'avvio.

\begin{mlbox}[Perché MAC «leggibili»]
Di norma Kathará assegna MAC casuali. Qui sono forzati: \texttt{01}, \texttt{02}, \texttt{03} per i pc; \texttt{a1}, \texttt{b1}, \texttt{b2}, \texttt{c1} per le interfacce dei router (la lettera indica la LAN). Così nelle catture Wireshark e nella cache ARP si capisce a colpo d'occhio chi sta parlando.
\end{mlbox}

\subsection{Il dispositivo wireshark}
\begin{intuizione}
È un container speciale che contiene Wireshark con un'interfaccia grafica raggiungibile \textbf{dal browser}. All'avvio è collegato solo all'host (\texttt{bridged}, interfaccia \texttt{eth0}); quando vuoi osservare una LAN, lo «colleghi» con \texttt{kathara lconfig}, come se attaccassi un cavo dal tuo PC di analisi allo switch di quella LAN. Non apre terminali (\texttt{num\_terms=0}): si usa solo dal browser su \texttt{localhost:3000}.
\end{intuizione}
\begin{attenzione}[Quale interfaccia sniffare]
Le interfacce aggiunte con \texttt{lconfig} sono numerate \textbf{in ordine dopo l'ultima esistente}. Nel lab: \texttt{eth0} = bridged (verso l'host), la prima LAN aggiunta (C) diventa \texttt{eth1}, la seconda (B) diventa \texttt{eth2}. Per questo le slide dicono di sniffare \texttt{eth1} su C e \texttt{eth2} su B.
\end{attenzione}

\subsection{I file .startup}
Sono script di shell: ogni riga è un comando eseguito dentro il dispositivo all'avvio, come se lo scrivessi tu nel suo terminale.

\begin{filecfg}{pc1.startup}
ip address add 195.11.14.5/24 dev eth0
ip route add default via 195.11.14.1
\end{filecfg}
\begin{filecfg}{pc2.startup}
ip address add 200.1.1.7/24 dev eth0
ip route add default via 200.1.1.1 dev eth0
\end{filecfg}
\begin{filecfg}{pc3.startup}
ip address add 200.1.1.3/24 dev eth0
ip route add default via 200.1.1.1 dev eth0
\end{filecfg}
\begin{tecnica}[Cosa fanno, riga per riga (host)]
\begin{enumerate}
  \item \inl{ip address add 195.11.14.5/24 dev eth0}: assegna a \texttt{eth0} l'indirizzo \texttt{195.11.14.5} con maschera \texttt{/24} (255.255.255.0). Effetto collaterale importante: il kernel aggiunge da solo la rotta verso la rete \texttt{195.11.14.0/24}, che diventa \textbf{direttamente connessa}.
  \item \inl{ip route add default via 195.11.14.1}: imposta il \textbf{default gateway}. Tutto ciò che non è destinato alla propria LAN va consegnato a \texttt{195.11.14.1}, cioè a \texttt{r1}. In pc2 e pc3 c'è anche \texttt{dev eth0}, che indica esplicitamente l'interfaccia di uscita; in pc1 manca perché il kernel la deduce da solo (il gateway sta su una rete di \texttt{eth0}).
\end{enumerate}
\end{tecnica}

\begin{filecfg}{r1.startup}
ip address add 195.11.14.1/24 dev eth0    °# LAN A: rete direttamente connessa°
ip address add 100.0.0.9/30 dev eth1      °# LAN B: rete direttamente connessa°
ip route add 200.1.1.0/24 via 100.0.0.10 dev eth1   °# LAN C: rotta statica°
\end{filecfg}
\begin{filecfg}{r2.startup  (ricostruito)}
ip address add 200.1.1.1/24 dev eth0
ip address add 100.0.0.10/30 dev eth1
ip route add 195.11.14.0/24 via 100.0.0.9 dev eth1
\end{filecfg}
\begin{tecnica}[Cosa fanno (router)]
\begin{itemize}
  \item Le prime due righe danno un indirizzo a ciascuna interfaccia: le reti A e B (per r1) diventano \textbf{direttamente connesse} e non serve scrivere rotte per raggiungerle.
  \item La terza riga aggiunge una \textbf{rotta statica} verso l'unica LAN che il router non vede direttamente: per r1 è C (\texttt{200.1.1.0/24}), raggiungibile consegnando i pacchetti a r2 (\texttt{100.0.0.10}) attraverso \texttt{eth1}.
  \item Le slide mostrano solo \texttt{r1.startup} e dicono «configurazione analoga per r2»: il file di r2 qui sopra è ricostruito dalla topologia ed è coerente con il piano di indirizzamento.
\end{itemize}
\end{tecnica}
\begin{attenzione}[Il forwarding]
Un router Linux inoltra i pacchetti tra le sue interfacce solo se è attivo il parametro \texttt{net.ipv4.ip\_forward}. Kathará lo attiva di default all'avvio dei dispositivi, ed è per questo che \texttt{r1} e \texttt{r2} funzionano da router senza comandi aggiuntivi (e, in linea di principio, anche un pc inoltrerebbe pacchetti se ne ricevesse da inoltrare). Se in un tuo lab un «router» non inoltra, controlla con \inl{sysctl net.ipv4.ip\_forward} (deve valere 1).
\end{attenzione}

\begin{appunti}{Riepilogo del capitolo}
Lab = cartella con \texttt{lab.conf} + \texttt{<dev>.startup} (+ \texttt{shared/}, cartelle per dispositivo). \texttt{dev[N]="CD/MAC"} collega \texttt{ethN} alla LAN \texttt{CD}. Opzioni: \texttt{image}, \texttt{ipv6}, \texttt{bridged}, \texttt{port}, \texttt{num\_terms}. Host: indirizzo + default gateway. Router: un indirizzo per interfaccia (reti dirette) + rotta statica verso la LAN remota. wireshark: \texttt{eth0} bridged, poi \texttt{eth1}, \texttt{eth2}\dots{} nell'ordine di \texttt{lconfig}.
\end{appunti}
